\section{Conclusion and Future Work}

This report presented LectureMind as an algorithmic framework for transforming lecture videos into structured and interactive learning resources. The framework begins with multimodal preprocessing, combines slide transitions and speech timing for hybrid chunking, converts chunks into explicit knowledge units, and supports grounded question answering through a three-mode retrieval architecture. The key idea is that lecture understanding should be hierarchical: raw video is first reorganized into coherent segments, then into structured knowledge, and finally into evidence-aware answers.

The evaluation suggests that this design is effective. Hybrid chunking outperforms single-signal segmentation, hierarchical knowledge representation supports targeted retrieval, and the three answering modes provide a useful balance between speed, quality, and reliability. In particular, Agentic RAG offers the strongest performance on complex questions, while Fast RAG remains an attractive conservative option for grounded responses.

Several limitations remain. The current formulation is centered on single-lecture understanding rather than full course-level reasoning. Knowledge extraction quality still depends on the clarity of speech and slide content. In addition, the agentic mode can become overconfident when evidence is incomplete, which means safety-oriented answer control remains an open problem.

Future work can extend the framework in three promising directions. First, cross-video retrieval could support course-level study, allowing concepts introduced in different lectures to be jointly retrieved and compared. Second, more adaptive chunking could learn boundary decisions from data rather than relying mainly on heuristic fusion of multimodal signals. Third, richer multimodal knowledge representations could better capture equations, diagrams, and other slide artifacts that are only partially reflected in speech. These directions would move lecture intelligence closer to a genuinely interactive educational assistant.
